\documentclass[utf8]{ctexart}
\usepackage{graphicx}
\usepackage{amsmath, amssymb, amsthm}
\usepackage{geometry}
\usepackage{enumitem}
\usepackage{xcolor}
\usepackage{mdframed}
\usepackage{tikz}

\usetikzlibrary{arrows.meta,positioning,calc}
\usepackage{booktabs}
\usepackage{array}

\geometry{margin=2cm}

\definecolor{critical}{RGB}{200,40,40}
\definecolor{nodeblue}{RGB}{52,120,200}
\definecolor{nodegray}{RGB}{215,228,242}
\definecolor{critnodebg}{RGB}{255,218,218}

\usetikzlibrary{graphs, positioning, arrows.meta}

\geometry{a4paper, margin=2.5cm}

% 定理环境
\newtheoremstyle{mystyle}
  {8pt}{8pt}{\kaishu}{0pt}{\bfseries}{}{1em}{}
\theoremstyle{mystyle}
\newtheorem{theorem}{定理}
\newtheorem{lemma}{引理}

% 好看的题目框
\newmdenv[
  linecolor=black!40,
  linewidth=1pt,
  backgroundcolor=gray!8,
  roundcorner=5pt,
  innertopmargin=8pt,
  innerbottommargin=8pt,
]{problembox}

\title{\textbf{离散数学2\quad 作业9}}
\date{\today}
\author{李子嘉 2024012325}
\begin{document}
\maketitle

\section*{题目1}

\begin{problembox}
  证明：在 $n \geq 4$ 的极大平面图中，每个顶点的度数至少为 $3$。
\end{problembox}
\textbf{解：} 
我们知道，极大平面图具备以下性质：
\begin{itemize}
  \item $ m = 3n - 6 $，其中 $n$ 是顶点数，$m$ 是边数。
  \item $ r = 2n - 4 $，其中 $r$ 是面数。
  \item $G$ 的每个域都是三角形，每条边都被两个不同的面共享，即 $2m = 3r$。
\end{itemize}
因此，每条边都被两个三角形共享，显而易见，每个顶点至少连接三条边，否则无法形成2个三角形。换言之，每个顶点的度数至少为 $3$。
严格地，我们假设存在一个顶点 $x$，$\deg(x) = 2$，设两个邻居为$u$和$v$，则$uvx$构成一个三角形。由于边$(u,x)$被两个三角形共享，那么必然另有一个顶点$w$使得$uwx$也是一个三角形，因此$w$与$x$有边相连，这与$\deg(x) = 2$矛盾。因此，$\deg(x) \geq 3$，得证。

\section*{题目2}
\begin{problembox}
设图 $G$ 是 $n$个点$m$条边的简单平面图，且$G$是自对偶图，证明：$m=2n-2$。
\end{problembox}
\textbf{解：}
由于 $G$ 是自对偶图，所以 $G$ 的面数 $r$ 等于顶点数 $n$，即 $r = n$。根据欧拉公式，我们有：
\[n - m + r = 2 \implies n - m + n = 2 \implies 2n - m = 2 \implies m = 2n - 2\]
得证。

\section*{题目3}
\begin{problembox}
假定一个带有 $m$ 条边和 $n$ 个顶点的连通简单平面图不包含长度为4或更短的回路，证明：若 $n \geq 4$，则 $m \leq \frac{5}{3}n - \frac{10}{3}$。
\end{problembox}
\textbf{解：}
由题可知，图 $G$ 中没有长度为4或更短的回路，因此每个域至少由5条边围成。因此我们有：
\[2m \geq 5r\]
其中 $r$ 是面数。结合欧拉公式 $n - m + r = 2$，我们可以得到：
\[n - m + \frac{2m}{5} = 2\]
整理得：
\[n - m + \frac{2m}{5} = 2 \implies n - \frac{3m}{5} = 2 \implies \frac{3m}{5} = n - 2 \implies m = \frac{5}{3}(n - 2) = \frac{5}{3}n - \frac{10}{3}\]
得证。

\section*{题目4}
\begin{problembox}
设简单连通图 $G$ 的顶点数是15，其中8个点的度数是4，6个点的度数是6，一个点的度数是8，证明：$G$ 不是平面图。
\end{problembox}
\textbf{解：}
根据题目，我们可以计算图 $G$ 的边数 $m$：
\[m = \frac{1}{2} \sum_{v \in V} \deg(v) = \frac{1}{2} (8 \cdot 4 + 6 \cdot 6 + 1 \cdot 8) = \frac{1}{2} (32 + 36 + 8) = \frac{1}{2} \cdot 76 = 38\]
根据欧拉公式，$n=15, m=38, r=25$.
由于原图无环或重边，每个域至少有3条边，我们关注到$3r=75$接近$2m=76$，做个简单的鸡兔同笼问题：
设有$x$个域由3条边围成，$y$个域由4条边围成（不可能有5条边围成，否则$71/3 < 24$），则：
\[\begin{cases}
3x + 4y = 76 \\
x + y = 25
\end{cases}\]
解得$x=24, y=1$，即有24个域由3条边围成，1个域由4条边围成。
假设原图是平面图，我们对其唯一的四边形面，加一条边连接其对角线，得到新图$G'$，则$G'$有15个顶点，39条边，26个面。由于$3r=2m=78, m=3n-6=39$，，因此$G'$是一个极大平面图。
而对于新图，总度数需为$6n-12=78$，但原图中已经有一个顶点的度数为8，其他顶点的度数不变，因此新图中至少有一个顶点的度数大于或等于8，这与极大平面图中每个顶点的度数至少为3且不超过5矛盾。因此，原图$G$不是平面图。
\end{document}